%!TEX root = prelim.tex



%\frame{
%\frametitle{Why is face recognition a good way to recognize people?}
%\begin{itemize}
%\item Non-contact
%\item Potentially faster than gait or voice recognition
%\item No special hardware needed for recognition (unlike Iris recognition)
%\item Attacks would probably be obvious to bystanders
%\item Recognized person need not do anything
%\end{itemize}
%}

%\frame{
%\frametitle{Importance for humans}
%\begin{itemize}
%\item 
%Humans seem to have evolved domain specific "wetware" for recognizing faces, called the fusiform face area.  \footnote{"Early (M170) activation of face-specific cortex by face-like objects", Hadjikhani N, et al., Feb 2009, Neuroreport 20: 403.}
%\includegraphics[width=0.9\textwidth]{images/fusiform_face_area.jpg}
%\item Face Blindness (Prosopagnosia) congenital in 2.5\% of population \\
%%\item Capgras syndrome
%\end{itemize}
%}

%\subsection{Review of Problem}
\frame{
\frametitle{Status of image based face recognition in industry}
\begin{itemize}
\item Barely works for entertainment purposes with a small number of users.
\begin{itemize}
\item iPhoto
\item Picassa
\item Windows Live Photo Gallery (Beta)
\end{itemize}
\item A failure for security applications
\begin{itemize}
\item High profile trial at Boston Logan Airport failed over half of the time.
\item Toshiba laptop face recognition can be broken by holding up a picture
\end{itemize}
\end{itemize}
}

%DREW: PICTURES

\frame{
\frametitle{Status of image based face recognition in academia}
\begin{itemize}
\item Classical algorithms, such as Eigenfaces, LDA, Nearest Neighbor, and Nearest Subspace perform poorly when confronted with test images taken outside of the lab.
\item Three factors dramatically affect the appearance of a face:
\begin{enumerate}
\item Illumination Variation: The appearance of the face depends very strongly on the illumination
\item Alignment Variation: The image depends on the orientation between the camera and subject 
\item Occlusion: Often parts of the face image are covered by other objects
\end{enumerate}
\end{itemize}
}


\frame{
\frametitle{A promising recent direction in face recognition research}
\begin{center}
\vspace{-.2in}
\includegraphics[width=.3\textwidth]{images/bouqet.pdf}\\
\vspace{-.2in}
\end{center}
\begin{itemize}
\item One very promising direction in face recognition research,
{\bf Sparse Representation based Classification (SRC)} [Wright et.\ al. PAMI 2009], achieved state-of-the-art results on Extended Yale B and AR databases.
\item This is a solid foundation, but it is incomplete:
	\begin{itemize}
	\item Assumed perfectly aligned images
	\item Assumed a sufficient number of training illuminations taken under different illuminations
\end{itemize}
\end{itemize}
}



\frame{
\frametitle{The Goal of the Proposed Research}
{\large The goal of my research is to build a complete and scalable face recognition system that can handle  {\bf illumination variation, alignment error, and realistic occlusions.}}\\
\vspace{.2in}
This will entail the following major contributions:
\begin{itemize}
\item Determine what {\bf training illuminations} are sufficient and build an {\bf acquisition system} capable of capturing them
\item Design algorithm for {\bf automatically aligning} test images to training images
\item Improve performance on {\bf contiguous occlusions}
\end{itemize}
}
